\renewcommand{\labelenumi}{\Roman{enumi}.}

\hypertarget{contents}{%
\subsection{CONTENTS}\label{contents}}

\hypertarget{introductory}{%
\subsubsection{INTRODUCTORY}\label{introductory}}

\begin{adjustwidth}{2em}{}
[The Problem](#Introductory)
\end{adjustwidth}

\hypertarget{book-1.-wages-and-capital}{%
\subsubsection{BOOK 1. WAGES AND
CAPITAL}\label{book-1.-wages-and-capital}}

\begin{enumerate}
\def\labelenumi{\arabic{enumi}.}
\tightlist
\item
  \protect\hyperlink{Book1Chapter1}{The current doctrine of wages---its
  insufficiency}
\item
  \protect\hyperlink{Book1Chapter2}{The meaning of the terms}
\item
  \protect\hyperlink{Book1Chapter3}{Wages not drawn from capital, but
  produced by the labor}
\item
  The maintenance of laborers not down from capital
\item
  The real functions of capital
\end{enumerate}

\hypertarget{book-2.-population-and-subsistence}{%
\subsubsection{BOOK 2. POPULATION AND
SUBSISTENCE}\label{book-2.-population-and-subsistence}}

\begin{enumerate}
\def\labelenumi{\arabic{enumi}.}
\tightlist
\item
  The Malthusian theory, its genesis and support
\item
  Inferences from facts
\item
  Inferences from analogy
\item
  Disproof of the Malthusian theory
\end{enumerate}

\hypertarget{book-3.-the-laws-of-distribution}{%
\subsubsection{BOOK 3. THE LAWS OF
DISTRIBUTION}\label{book-3.-the-laws-of-distribution}}

\begin{enumerate}
\def\labelenumi{\arabic{enumi}.}
\tightlist
\item
  The inquiry narrowed to the laws of distribution necessary relation of
  these laws
\item
  Rent and the law of rent
\item
  Interest and the cause of interest
\item
  Of spurious capital and of profits often mistaken forinterest
\item
  The law of interest
\item
  Wages and the law of wages
\item
  Correlation and co-ordination of these laws
\item
  The statics of the problem thus explained
\end{enumerate}

\hypertarget{book-4.-effect-of-material-progress-upon-the-distribution-of-wealth}{%
\subsubsection{BOOK 4. EFFECT OF MATERIAL PROGRESS UPON THE DISTRIBUTION
OF
WEALTH}\label{book-4.-effect-of-material-progress-upon-the-distribution-of-wealth}}

\begin{enumerate}
\def\labelenumi{\arabic{enumi}.}
\tightlist
\item
  The dynamics of the problem yet to seek.
\item
  Effect of increase of population upon the distribution of wealth.
\item
  Effect of improvements in the arts upon the distribution of wealth
\item
  Effect of the expectation raised by material progress.
\end{enumerate}

\hypertarget{book-5.-the-problem-solved}{%
\subsubsection{BOOK 5. THE PROBLEM
SOLVED}\label{book-5.-the-problem-solved}}

\begin{enumerate}
\def\labelenumi{\arabic{enumi}.}
\tightlist
\item
  The primary cause of recurring paroxysms of industrial depression
\item
  The persistence of poverty amid advancing wealth.
\end{enumerate}

\hypertarget{book-6.-the-remedy}{%
\subsubsection{BOOK 6. THE REMEDY}\label{book-6.-the-remedy}}

\begin{enumerate}
\def\labelenumi{\arabic{enumi}.}
\tightlist
\item
  Insufficiency of remedies currently advocated
\item
  The true remedy
\end{enumerate}

\hypertarget{book-7.-justice-of-the-remedy}{%
\subsubsection{BOOK 7. JUSTICE OF THE
REMEDY}\label{book-7.-justice-of-the-remedy}}

\begin{enumerate}
\def\labelenumi{\arabic{enumi}.}
\tightlist
\item
  Injustice of private property in land
\item
  Enslavement of laborers the ultimate result of private property in
  land
\item
  Claim of land owners to compensation
\item
  Property in land historically considered
\item
  Property in land in the UnIted States
\end{enumerate}

\hypertarget{book-8.-application-of-the-remedy}{%
\subsubsection{BOOK 8. APPLICATION OF THE
REMEDY}\label{book-8.-application-of-the-remedy}}

\begin{enumerate}
\def\labelenumi{\arabic{enumi}.}
\tightlist
\item
  Private property in land inconsistent with the best use of land.
\item
  How equal rights to the land may be asserted and secured.
\item
  The proposition tried by the canons of taxation.
\item
  Indorsements and objections
\end{enumerate}

\hypertarget{book-ix.-effects-of-the-remedy}{%
\subsubsection{BOOK IX. EFFECTS OF THE
REMEDY}\label{book-ix.-effects-of-the-remedy}}

\begin{enumerate}
\def\labelenumi{\arabic{enumi}.}
\tightlist
\item
  Of the effect upon the production of wealth.
\item
  Of the effect upon distribution and thence upon production
\item
  Of the effect upon individuals and classes
\item
  Of the changes that would be wrought in social organization and social
  life
\end{enumerate}

\hypertarget{book-x.-the-law-of-human-progress}{%
\subsubsection{BOOK X. THE LAW OF HUMAN
PROGRESS}\label{book-x.-the-law-of-human-progress}}

\begin{enumerate}
\def\labelenumi{\arabic{enumi}.}
\tightlist
\item
  The current theory of human progress---its insufficiency
\item
  Differences in civilization---to what due
\item
  The law of human progress
\item
  How modern civilization may decline
\item
  The central truth
\end{enumerate}

\hypertarget{conclusion}{%
\subsubsection{CONCLUSION}\label{conclusion}}

The problem of individual life

\newpage

\hypertarget{Introductory}{%
\section{INTRODUCTORY}\label{Introductory}}

\hypertarget{the-problem}{%
\subsection{THE PROBLEM}\label{the-problem}}

The present century has been marked by a prodigious increase in
wealth-producing power. The utilization of steam and electricity, the
introduction of improved processes and labor-saving machinery, the
greater subdivision and grander scale of production, the wonderful
facilitation of exchanges, have multiplied enormously the effectiveness
of labor.

At the beginning of this marvelous era it was natural to expect, and it
was expected, that labor-saving inventions would lighten the toil and
improve the condition of the laborer; that the enormous increase in the
power of producing wealth would make real poverty a thing of the past.
Could a man of the last century---a Franklin or a Priestley---have seen,
in a vision of the future, the steamship taking the place of the sailing
vessel, the railroad train of the wagon, the reaping machine of the
scythe, the threshing machine of the flail; could he have heard the
throb of the engines that in obedience to human will, and for the
satisfaction of human desire, exert a power greater than that of all the
men and all the beasts of burden of the earth combined; could he have
seen the forest tree transformed into finished lumber---into doors,
sashes, blinds, boxes or barrels, with hardly the touch of a human hand;
the great workshops where boots and shoes are turned out by the case
with less labor than the old-fashioned cobbler could have put on a sole;
the factories where, under the eye of a girl, cotton becomes cloth
faster than hundreds of stalwart weavers could have turned it out with
their hand looms; could he have seen steam hammers shaping mammoth
shafts and mighty anchors, and delicate machinery making tiny watches;
the diamond drill cutting through the heart of the rocks, and coal oil
sparing the whale; could he have realized the enormous saving of labor
resulting from improved facilities of exchange and communication---sheep
killed in Australia eaten fresh in England, and the order given by the
London banker in the afternoon executed in San Francisco in the morning
of the same day; could he have conceived of the hundred thousand
improvements which these only suggest, what would he have inferred as to
the social condition of mankind?

It would not have seemed like an inference; further than the vision went
it would have seemed as though he saw; and his heart would have leaped
and his nerves would have thrilled, as one who from a height beholds
just ahead of the thirst-stricken caravan the living gleam of rustling
woods and the glint of laughing waters. Plainly, in the sight of the
imagination, he would have beheld these new forces elevating society
from its very foundations, lifting the very poorest above the
possibility of want, exempting the very lowest from anxiety for the
material needs of life; he would have seen these slaves of the lamp of
knowledge taking on themselves the traditional curse, these muscles of
iron and sinews of steel making the poorest laborer's life a holiday, in
which every high quality and noble impulse could have scope to grow.

And out of these bounteous material conditions he would have seen
arising, as necessary sequences, moral conditions realizing the golden
age of which mankind have always dreamed. Youth no longer stunted and
starved; age no longer harried by avarice; the child at play with the
tiger; the man with the muck-rake drinking in the glory of the stars.
Foul things fled, fierce things tame; discord turned to harmony! For how
could there be greed where all had enough? How could the vice, the
crime, the ignorance, the brutality, that spring from poverty and the
fear of poverty, exist where poverty had vanished? Who should crouch
where all were freemen; who oppress where all were peers?

More or less vague or clear, these have been the hopes, these the dreams
born of the improvements which give this wonderful century its
preeminence. They have sunk so deeply into the popular mind as radically
to change the currents of thought, to recast creeds and displace the
most fundamental conceptions. The haunting visions of higher
possibilities have not merely gathered splendor and vividness, but their
direction has changed ---instead of seeing behind the faint tinges of an
expiring sunset, all the glory of the daybreak has decked the skies
before.

It is true that disappointment has followed disappointment, and that
discovery upon discovery, and invention after invention, have neither
lessened the toil of those who most need respite, nor brought plenty to
the poor. But there have been so many things to which it seemed this
failure could be laid, that up to our time the new faith has hardly
weakened. We have better appreciated the difficulties to be overcome;
but not the less trusted that the tendency of the times was to overcome
them.

Now, however, we are coming into collision with facts which there can be
no mistaking. From all parts of the civilized world come complaints of
industrial depression; of labor condemned to involuntary idleness; of
capital massed and wasting; of pecuniary distress among business men; of
want and suffering and anxiety among the working classes. All the dull,
deadening pain, all the keen, maddening anguish, that to great masses of
men are involved in the words ``hard times,'' afflict the world to-day.
This state of things, common to communities differing so widely in
situation, in political institutions, in fiscal and financial systems,
in density of population and in social organization, can hardly be
accounted for by local causes. There is distress where large standing
armies are maintained, but there is also distress where the standing
armies are nominal; there is distress where protective tariffs stupidly
and wastefully hamper trade, but there is also distress where trade is
nearly free; there is distress where autocratic government yet prevails,
but there is also distress where political power is wholly in the hands
of the people; in countries where paper is money, and in countries where
gold and silver are the only currency. Evidently, beneath all such
things as these, we must infer a common cause.

That there is a common cause, and that it is either what we call
material progress or something closely connected with material progress,
becomes more than an inference when it is noted that the phenomena we
class together and speak of as industrial depression are but
intensifications of phenomena which always accompany material progress,
and which show themselves more clearly and strongly as material progress
goes on. Where the conditions to which material progress everywhere
tends are most fully realized---that is to say, where population is
densest, wealth greatest, and the machinery of production and exchange
most highly developed---we find the deepest poverty, the sharpest
struggle for existence, and the most of enforced idleness.

It is to the newer countries---that is, to the countries where material
progress is yet in its earlier stages---that laborers emigrate in search
of higher wages, and capital flows in search of higher interest. It is
in the older countries---that is to say, the countries where material
progress has reached later stages---that widespread destitution is found
in the midst of the greatest abundance. Go into one of the new
communities where Anglo-Saxon vigor is just beginning the race of
progress; where the machinery of production and exchange is yet rude and
inefficient; where the increment of wealth is not yet great enough to
enable any class to live in ease and luxury; where the best house is but
a cabin of logs or a cloth and paper shanty, and the richest man is
forced to daily work---and though you will find an absence of wealth and
all its concomitants, you will find no beggars. There is no luxury, but
there is no destitution. No one makes an easy living, nor a very good
living; but everyone can make a living, and no one able and willing to
work is oppressed by the fear of want.

But just as such a community realizes the conditions which all civilized
communities are striving for, and advances in the scale of material
progress---just as closer settlement and a more intimate connection with
the rest of the world, and greater utilization of labor-saving
machinery, make possible greater economies in production and exchange,
and wealth in consequence increases, not merely in the aggregate, but in
proportion to population ---so does poverty take a darker aspect. Some
get an infinitely better and easier living, but others find it hard to
get a living at all. The ``tramp'' comes with the locomotive, and
almshouses and prisons are as surely the marks of ``material progress''
as are costly dwellings, rich warehouses, and magnificent churches. Upon
streets lighted with gas and patrolled by uniformed policemen, beggars
wait for the passer-by, and in the shadow of college, and library, and
museum, are gathering the more hideous Huns and fiercer Vandals of whom
Macaulay prophesied.

This fact---the great fact that poverty and all its concomitants show
themselves in communities just as they develop into the conditions
toward which material progress tends---proves that the social
difficulties existing wherever a certain stage of progress has been
reached, do not arise from local circumstances, but are, in some way or
another, engendered by progress itself.

And, unpleasant as it may be to admit it, it is at last becoming evident
that the enormous increase in productive power which has marked the
present century and is still going on with accelerating ratio, has no
tendency to extirpate poverty or to lighten the burdens of those
compelled to toil. It simply widens the gulf between Dives and Lazarus,
and makes the struggle for existence more intense. The march of
invention has clothed mankind with powers of which a century ago the
boldest imagination could not have dreamed. But in factories where
labor-saving machinery has reached its most wonderful development,
little children are at work; wherever the new forces are anything like
fully utilized, large classes are maintained by charity or live on the
verge of recourse to it; amid the greatest accumulations of wealth, men
die of starvation, and puny infants suckle dry breasts; while everywhere
the greed of gain, the worship of wealth, shows the force of the fear of
want. The promised land flies before us like the mirage. The fruits of
the tree of knowledge turn as we grasp them to apples of Sodom that
crumble at the touch.

It is true that wealth has been greatly increased, and that the average
of comfort, leisure, and refinement has been raised; but these gains are
not general. In them the lowest class do not share.* I do not mean that
the condition of the lowest class has nowhere nor in anything been
improved; but that there is nowhere any improvement which can be
credited to increased productive power. I mean that the tendency of what
we call material progress is in nowise to improve the condition of the
lowest class in the essentials of healthy, happy human life. Nay, more,
that it is still further to depress the condition of the lowest class.
The new forces, elevating in their nature though they be, do not act
upon the social fabric from underneath, as was for a long time hoped and
believed, but strike it at a point intermediate between top and bottom.
It is as though an immense wedge were being forced, not underneath
society, but through society. Those who are above the point of
separation are elevated, but those who are below are crushed down.

\begin{itemize}
\tightlist
\item
  It is true that the poorest may now in certain ways enjoy what the
  richest a century ago could not have commanded, but this does not show
  improvement of condition so long as the ability to obtain the
  necessaries of life is not increased. The beggar in a great city may
  enjoy many things from which the backwoods farmer is debarred, but
  that does not prove the condition of the city beggar better than that
  of the independent farmer.
\end{itemize}

This depressing effect is not generally realized, for it is not apparent
where there has long existed a class just able to live. Where the lowest
class barely lives, as has been the case for a long time in many parts
of Europe, it is impossible for it to get any lower, for the next lowest
step is out of existence, and no tendency to further depression can
readily show itself. But in the progress of new settlements to the
conditions of older communities it may clearly be seen that material
progress does not merely fail to relieve poverty---it actually produces
it. In the United States it is clear that squalor and misery, and the
vices and crimes that spring from them, everywhere increase as the
village grows to the city, and the march of development brings the
advantages of the improved methods of production and exchange. It is in
the older and richer sections of the Union that pauperism and distress
among the working classes are becoming most painfully apparent. If there
is less deep poverty in San Francisco than in New York, is it not
because San Francisco is yet behind New York in all that both cities are
striving for? When San Francisco reaches the point where New York now
is, who can doubt that there will also be ragged and barefooted children
on her streets?

This association of poverty with progress is the great enigma of our
times. It is the central fact from which spring industrial, social, and
political difficulties that perplex the world, and with which
statesmanship and philanthropy and education grapple in vain. From it
come the clouds that overhang the future of the most progressive and
self-reliant nations. It is the riddle which the Sphinx of Fate puts to
our civilization, and which not to answer is to be destroyed. So long as
all the increased wealth which modern progress brings goes but to build
up great fortunes, to increase luxury and make sharper the contrast
between the House of Have and the House of Want, progress is not real
and cannot be permanent. The reaction must come. The tower leans from
its foundations, and every new story but hastens the final catastrophe.
To educate men who must be condemned to poverty, is but to make them
restive; to base on a state of most glaring social inequality political
institutions under which men are theoretically equal, is to stand a
pyramid on its apex.

All-important as this question is, pressing itself from every quarter
painfully upon attention, it has not yet received a solution which
accounts for all the facts and points to any clear and simple remedy.
This is shown by the widely varying attempts to account for the
prevailing depression. They exhibit not merely a divergence between
vulgar notions and scientific theories, but also show that the
concurrence which should exist between those who avow the same general
theories breaks up upon practical questions into an anarchy of opinion.
Upon high economic authority we have been told that the prevailing
depression is due to over-consumption; upon equally high authority, that
it is due to over-production; while the wastes of war, the extension of
railroads, the attempts of workmen to keep up wages, the demonetization
of silver, the issues of paper money, the increase of labor-saving
machinery, the opening of shorter avenues to trade, etc., are separately
pointed out as the cause, by writers of reputation.

And while professors thus disagree, the ideas that there is a necessary
conflict between capital and labor, that machinery is an evil, that
competition must be restrained and interest abolished, that wealth may
be created by the issue of money, that it is the duty of government to
furnish capital or to furnish work, are rapidly making way among the
great body of the people, who keenly feel a hurt and are sharply
conscious of a wrong. Such ideas, which bring great masses of men, the
repositories of ultimate political power, under the leadership of
charlatans and demagogues, are fraught with danger; but they cannot be
successfully combated until political economy shall give some answer to
the great question which shall be consistent with all her teachings, and
which shall commend itself to the perceptions of the great masses of
men.

It must be within the province of political economy to give such an
answer. For political economy is not a set of dogmas. It is the
explanation of a certain set of facts. It is the science which, in the
sequence of certain phenomena, seeks to trace mutual relations and to
identify cause and effect, just as the physical sciences seek to do in
other sets of phenomena. It lays its foundations upon firm ground. The
premises from which it makes its deductions are truths which have the
highest sanction; axioms which we all recognize; upon which we safely
base the reasoning and actions of everyday life, and which may be
reduced to the metaphysical expression of the physical law that motion
seeks the line of least resistance---viz., that men seek to gratify
their desires with the least exertion. Proceeding from a basis thus
assured, its processes, which consist simply in identification and
separation, have the same certainty. In this sense it is as exact a
science as geometry, which, from similar truths relative to space,
obtains its conclusions by similar means, and its conclusions when valid
should be as self-apparent. And although in the domain of political
economy we cannot test our theories by artificially produced
combinations or conditions, as may be done in some of the other
sciences, yet we can apply tests no less conclusive, by comparing
societies in which different conditions exist, or by, in imagination,
separating, combining, adding or eliminating forces or factors of known
direction.

I propose in the following pages to attempt to solve by the methods of
political economy the great problem I have outlined. I propose to seek
the law which associates poverty with progress, and increases want with
advancing wealth; and I believe that in the explanation of this paradox
we shall find the explanation of those recurring seasons of industrial
and commercial paralysis which, viewed independently of their relations
to more general phenomena, seem so inexplicable. Properly commenced and
carefully pursued, such an investigation must yield a conclusion that
will stand every test, and as truth, will correlate with all other
truth. For in the sequence of phenomena there is no accident. Every
effect has a cause, and every fact implies a preceding fact.

That political economy, as at present taught, does not explain the
persistence of poverty amid advancing wealth in a manner which accords
with the deep-seated perceptions of men; that the unquestionable truths
which it does teach are unrelated and disjointed; that it has failed to
make the progress in popular thought that truth, even when unpleasant,
must make; that, on the contrary, after a century of cultivation, during
which it has engrossed the attention of some of the most subtle and
powerful intellects, it should be spurned by the statesman, scouted by
the masses, and relegated in the opinion of many educated and thinking
men to the rank of a pseudo-science in which nothing is fixed or can be
fixed---must, it seems to me, be due not to any inability of the science
when properly pursued, but to some false step in its premises, or
overlooked factor in its estimates. And as such mistakes are generally
concealed by the respect paid to authority, I propose in this inquiry to
take nothing for granted, but to bring even accepted theories to the
test of first principles, and should they not stand the test, freshly to
interrogate facts in the endeavor to discover their law.

I propose to beg no question, to shrink from no conclusion, but to
follow truth wherever it may lead. Upon us is the responsibility of
seeking the law, for in the very heart of our civilization to-day women
faint and little children moan. But what that law may prove to be is not
our affair. If the conclusions that we reach run counter to our
prejudices, let us not flinch; if they challenge institutions that have
long been deemed wise and natural, let us not turn back.

\newpage

\hypertarget{Book1Chapter1}{%
\section{CHAPTER I}\label{Book1Chapter1}}

\hypertarget{the-current-doctrine-of-wagesits-insufficiency}{%
\subsection{THE CURRENT DOCTRINE OF WAGES---ITS
INSUFFICIENCY}\label{the-current-doctrine-of-wagesits-insufficiency}}

Reducing to its most compact form the problem we have set out to
investigate, let us examine, step by step, the explanation which
political economy, as now accepted by the best authority, gives of it.

The cause which produces poverty in the midst of advancing wealth is
evidently the cause which exhibits itself in the tendency, everywhere
recognized, of wages to a minimum. Let us, therefore, put our inquiry
into this compact form:

Why, in spite of increase in productive power, do wages tend to a
minimum which will give but a bare living?

The answer of the current political economy is, that wages are fixed by
the ratio between the number of laborers and the amount of capital
devoted to the employment of labor, and constantly tend to the lowest
amount on which laborers will consent to live and reproduce, because the
increase in the number of laborers tends naturally to follow and
overtake any increase in capital. The increase of the divisor being thus
held in check only by the possibilities of the quotient, the dividend
may be increased to infinity without greater result.

In current thought this doctrine holds all but undisputed sway. It bears
the indorsement of the very highest names among the cultivators of
political economy, and though there have been attacks upon it, they are
generally more formal than real.* It is assumed by Buckle as the basis
of his generalizations of universal history. It is taught in all, or
nearly all, the great English and American universities, and is laid
down in textbooks which aim at leading the masses to reason correctly
upon practical affairs, while it seems to harmonize with the new
philosophy, which, having in a few years all but conquered the
scientific world, is now rapidly permeating the general mind.

\begin{itemize}
\tightlist
\item
  This seems to me true of Mr.~Thornton's objections, for while he
  denies the existence of a predetermined wage fund, consisting of a
  portion of capital set apart for the purchase of labor, he yet holds
  (which is the essential thing) that wages are drawn from capital, and
  that increase or decrease of capital is increase or decrease of the
  fund available for the payment of wages. The most vital attack upon
  the wage fund doctrine of which I know is that of Professor Francis A.
  Walker (The Wages Question: New York, 1876), yet he admits that wages
  are in large part advanced from capital---which, so far as it goes, is
  all that the stanchest supporter of the wage fund theory could
  claim---while he fully accepts the Malthusian theory. Thus his
  practical conclusions in nowise differ from those reached by
  expounders of the current theory.
\end{itemize}

Thus entrenched in the upper regions of thought, it is in cruder form
even more firmly rooted in what may be styled the lower. What gives to
the fallacies of protection such a tenacious hold, in spite of their
evident inconsistencies and absurdities, is the idea that the sum to be
distributed in wages is in each community a fixed one, which the
competition of ``foreign labor'' must still further subdivide. The same
idea underlies most of the theories which aim at the abolition of
interest and the restriction of competition, as the means whereby the
share of the laborer in the general wealth can be increased; and it
crops out in every direction among those who are not thoughtful enough
to have any theories, as may be seen in the columns of newspapers and
the debates of legislative bodies.

And yet, widely accepted and deeply rooted as it is, it seems to me that
this theory does not tally with obvious facts. For, if wages depend upon
the ratio between the amount of labor seeking employment and the amount
of capital devoted to its employment, the relative scarcity or abundance
of one factor must mean the relative abundance or scarcity of the other.
Thus, capital must be relatively abundant where wages are high, and
relatively scarce where wages are low. Now, as the capital used in
paying wages must largely consist of the capital constantly seeking
investment, the current rate of interest must be the measure of its
relative abundance or scarcity. So, if it be true that wages depend upon
the ratio between the amount of labor seeking employment and the capital
devoted to its employment, then high wages, the mark of the relative
scarcity of labor, must be accompanied by low interest, the mark of the
relative abundance of capital, and reversely, low wages must be
accompanied by high interest.

This is not the fact, but the contrary. Eliminating from interest the
element of insurance, and regarding only interest proper, or the return
for the use of capital, is it not a general truth that interest is high
where and when wages are high, and low where and when wages are low?
Both wages and interest have been higher in the United States than in
England, in the Pacific than in the Atlantic States. Is it not a
notorious fact that where labor flows for higher wages, capital also
flows for higher interest? Is it not true that wherever there has been a
general rise or fall in wages there has been at the same time a similar
rise or fall in interest? In California, for instance, when wages were
higher than anywhere else in the world, so also was interest higher.
Wages and interest have in California gone down together. When common
wages were \$5 a day, the ordinary bank rate of interest was twenty-four
per cent. per annum. Now that common wages are \$2 or \$2.50 a day, the
ordinary bank rate is from ten to twelve per cent.

Now, this broad, general fact, that wages are higher in new countries,
where capital is relatively scarce, than in old countries, where capital
is relatively abundant, is too glaring to be ignored. And although very
lightly touched upon, it is noticed by the expounders of the current
political economy. The manner in which it is noticed proves what I say,
that it is utterly inconsistent with the accepted theory of wages. For
in explaining it such writers as Mill, Fawcett, and Price virtually give
up the theory of wages upon which, in the same treatises, they formally
insist. Though they declare that wages are fixed by the ratio between
capital and laborers, they explain the higher wages and interest of new
countries by the greater relative production of wealth. I shall
hereafter show that this is not the fact, but that, on the contrary, the
production of wealth is relatively larger in old and densely populated
countries than in new and sparsely populated countries. But at present I
merely wish to point out the inconsistency. For to say that the higher
wages of new countries are due to greater proportionate production, is
clearly to make the ratio with production, and not the ratio with
capital, the determinator of wages.

Though this inconsistency does not seem to have been perceived by the
class of writers to whom I refer, it has been noticed by one of the most
logical of the expounders of the current political economy. Professor
Cairnes* endeavors in a very ingenious way to reconcile the fact with
the theory, by assuming that in new countries, where industry is
generally directed to the production of food and what in manufactures is
called ``raw material,'' a much larger proportion of the capital used in
production is devoted to the payment of wages than in older countries
where a greater part must be expended in machinery and material, and
thus, in the new country, though capital is scarcer, and interest is
higher, the amount determined to the payment of wages is really larger,
and wages are also higher. For instance, of \$100,000 devoted in an old
country to manufactures, \$80,000 would probably be expended for
buildings, machinery and the purchase of materials, leaving but \$20,000
to be paid out in wages; whereas in a new country, of \$30,000 devoted
to agriculture, etc., not more than \$5,000 would be required for tools,
etc., leaving \$25,000 to be distributed in wages. In this way it is
explained that the wage fund may be comparatively large where capital is
comparatively scarce, and high wages and high interest accompany each
other.

\begin{itemize}
\tightlist
\item
  Some Leading Principles of Political Economy Newly Expounded, Chapter
  1, Part 2.
\end{itemize}

In what follows I think I shall be able to show that this explanation is
based upon a total misapprehension of the relations of labor to
capital---a fundamental error as to the fund from which wages are drawn;
but at present it is necessary only to point out that the connection in
the fluctuation of wages and interest in the same countries and in the
same branches of industry cannot thus be explained. In those
alternations known as ``good times'' and ``hard times'' a brisk demand
for labor and good wages is always accompanied by a brisk demand for
capital and stiff rates of interest. While, when laborers cannot find
employment and wages droop, there is always an accumulation of capital
seeking investment at low rates.* The present depression has been no
less marked by want of employment and distress among the working classes
than by the accumulation of unemployed capital in all the great centers,
and by nominal rates of interest on undoubted security. Thus, under
conditions which admit of no explanation consistent with the current
theory, do we find high interest coinciding with high wages, and low
interest with low wages---capital seemingly scarce when labor is scarce,
and abundant when labor is abundant.

\begin{itemize}
\tightlist
\item
  Times of commercial panic are marked by high rates of discount, but
  this is evidently not a high rate of interest, properly so called, but
  a high rate of insurance against risk.
\end{itemize}

All these well known facts, which coincide with each other, point to a
relation between wages and interest, but it is to a relation of
conjunction, not of opposition. Evidently they are utterly inconsistent
with the theory that wages are determined by the ratio between labor and
capital, or any part of capital.

How, then, it will be asked, could such a theory arise? How is it that
it has been accepted by a succession of economists, from the time of
Adam Smith to the present day?

If we examine the reasoning by which in current treatises this theory of
wages is supported, we see at once that it is not an induction from
observed facts, but a deduction from a previously assumed theory, viz.,
that wages are drawn from capital. It being assumed that capital is the
source of wages, it necessarily follows that the gross amount of wages
must be limited by the amount of capital devoted to the employment of
labor, and hence that the amount individual laborers can receive must be
determined by the ratio between their number and the amount of capital
existing for their recompense.* This reasoning is valid, but the
conclusion, as we have seen, does not correspond with the facts. The
fault, therefore, must be in the premises. Let us see.

\begin{itemize}
\tightlist
\item
  For instance McCulloch (Note VI to Wealth of Nations) says: ``That
  portion of the capital or wealth of a country which the employers of
  labor intend to or are willing to pay out in the purchase of labor,
  may be much larger at one time than another. But whatever may be its
  absolute magnitude, it obviously forms the only source from which any
  portion of the wages of labor can be derived. No other fund is in
  existence from which the laborer, as such, can draw a single shilling.
  And hence it follows that the average rate of wages, or the share of
  the national capital appropriated to the employment of labor falling,
  at an average, to each laborer, must entirely depend on its amount as
  compared with the number of those amongst whom it has to be divided.''
  Similar citations might be made from all the standard economists.
\end{itemize}

I am aware that the theorem that wages are drawn from capital is one of
the most fundamental and apparently best settled of current political
economy, and that it has been accepted as axiomatic by all the great
thinkers who have devoted their powers to the elucidation of the
science. Nevertheless, I think it can be demonstrated to be a
fundamental error---the fruitful parent of a long series of errors,
which vitiate most important practical conclusions. This demonstration I
am about to attempt. It is necessary that it should be clear and
conclusive, for a doctrine upon which so much important reasoning is
based, which is supported by such a weight of authority, which is so
plausible in itself, and is so liable to recur in different forms,
cannot be safely brushed aside in a paragraph.

The proposition I shall endeavor to prove, is: That wages, instead of
being drawn from capital, are in reality drawn from the product of the
labor for which they are paid.*

\begin{itemize}
\tightlist
\item
  We are speaking of labor expended in production, to which it is best
  for the sake of simplicity to confine the inquiry. Any question which
  may arise in the reader's mind as to wages for unproductive services
  had best therefore be deferred.
\end{itemize}

Now, inasmuch as the current theory that wages are drawn from capital
also holds that capital is reimbursed from production, this at first
glance may seem a distinction without a difference---a mere change in
terminology, to discuss which would be but to add to those unprofitable
disputes that render so much that has been written upon
politico-economic subjects as barren and worthless as the controversies
of the various learned societies about the true reading of the
inscription on the stone that Mr.~Pickwick found. But that it is much
more than a formal distinction will be apparent when it is considered
that upon the difference between the two propositions are built up all
the current theories as to the relations of capital and labor; that from
it are, deduced doctrines that, themselves regarded as axiomatic, bound,
direct, and govern the ablest minds in the discussion of the most
momentous questions. For, upon the assumption that wages are drawn
directly from capital, and not from the product of the labor, is based,
not only the doctrine that wages depend upon the ratio between capital,
and labor, but the doctrine that industry is limited by capital, that
capital must be accumulated before labor is employed, and labor cannot
be employed except as capital is accumulated; the doctrine that every
increase of capital gives or is capable of giving additional employment
to industry; the doctrine that the conversion of circulating capital
into fixed capital lessens the fund applicable to the maintenance of
labor; the doctrine that more laborers can be employed at low than at
high wages; the doctrine that capital applied to agriculture will
maintain more laborers than if applied to manufactures; the doctrine
that profits are high or low as wages are low or high, or that they
depend upon the cost of the subsistence of laborers; together with such
paradoxes as that a demand for commodities is not a demand for labor, or
that certain commodities may be increased in cost by a reduction in
wages or diminished in cost by an increase in wages.

In short, all the teachings of the current political economy, in the
widest and most important part of its domain, are based more or less
directly upon the assumption that labor is maintained and paid out of
existing capital before the product which constitutes the ultimate
object is secured. If it be shown that this is an error, and that on the
contrary the maintenance and payment of labor do not even temporarily
trench on capital, but are directly drawn from the product of the labor,
then all this vast superstructure is left without support and must fall.
And so likewise must fall the vulgar theories which also have their base
in the belief that the sum to be distributed in wages is a fixed one,
the individual shares in which must necessarily be decreased by an
increase in the number of laborers.

The difference between the current theory and the one I advance is, in
fact, similar to that between the mercantile theory of international
exchanges and that with which Adam Smith supplanted it. Between the
theory that commerce is the exchange of commodities for money, and the
theory that it is the exchange of commodities for commodities, there may
seem no real difference when it is remembered that the adherents of the
mercantile theory did not assume that money had any other use than as it
could be exchanged for commodities. Yet, in the practical application of
these two theories, there arises all the difference between rigid
governmental protection and free trade.

If I have said enough to show the reader the ultimate importance of the
reasoning through which I am about to ask him to follow me, it will not
be necessary to apologize in advance either for simplicity or prolixity.
In arraigning a doctrine of such importance---a doctrine supported by
such a weight of authority, it is necessary to be both clear and
thorough.

Were it not for this I should be tempted to dismiss with a sentence the
assumption that wages are drawn from capital. For all the vast
superstructure which the current political economy builds upon this
doctrine is in truth based upon a foundation which has been merely taken
for granted, without the slightest attempt to distinguish the apparent
from the real. Because wages are generally paid in money, and in many of
the operations of production are paid before the product is fully
completed, or can be utilized, it is inferred that wages are drawn from
pre-existing capital, and, therefore, that industry is limited by
capital---that is to say that labor cannot be employed until capital has
been accumulated, and can only be employed to the extent that capital
has been accumulated. Yet in the very treatises in which the limitation
of industry by capital is laid down without reservation and made the
basis for the most important reasonings and elaborate theories, we are
told that capital is stored up or accumulated labor---``that part of
wealth which is saved to assist future production.'' If we substitute
for the word ``capital'' this definition of the word, the proposition
carries its own refutation, for that labor cannot be employed until the
results of labor are saved becomes too absurd for discussion.

Should we, however, with this reductio ad absurdum, attempt to close the
argument, we should probably be met with the explanation, not that the
first laborers were supplied by Providence with the capital necessary to
set them to work, but that the proposition merely refers to a state of
society in which production has become a complex operation.

But the fundamental truth, that in all economic reasoning must be firmly
grasped, and never let go, is that society in its most highly developed
form is but an elaboration of society in its rudest beginnings, and that
principles obvious in the simpler relations of men are merely disguised
and not abrogated or reversed by the more intricate relations that
result from the division of labor and the use of complex tools and
methods. The steam grist mill, with its complicated machinery exhibiting
every diversity of motion, is simply what the rude stone mortar dug up
from an ancient river bed was in its day---an instrument for grinding
corn. And every man engaged in it, whether tossing wood into the
furnace, running the engine, dressing stones, printing sacks or keeping
books, is really devoting his labor to the same purpose that the
prehistoric savage did when he used his mortar---the preparation of
grain for human food.

And so, if we reduce to their lowest terms all the complex operations of
modern production, we see that each individual who takes part in this
infinitely subdivided and intricate network of production and exchange
is really doing what the primeval man did when he climbed the trees for
fruit or followed the receding tide for shell-fish---endeavoring to
obtain from nature by the exertion of his powers the satisfaction of his
desires. If we keep this firmly in mind, if we look upon production as a
whole---as the co-operation of all embraced in any of its great groups
to satisfy the various desires of each, we plainly see that the reward
each obtains for his exertions comes as truly and as directly from
nature as the result of that exertion, as did that of the first man.

To illustrate: In the simplest state of which we can conceive, each man
digs his own bait and catches his own fish. The advantages of the
division of labor soon become apparent, and one digs bait while the
others fish. Yet evidently the one who digs bait is in reality doing as
much toward the catching of fish as any of those who actually take the
fish. So when the advantages of canoes are discovered, and instead of
all going a-fishing, one stays behind and makes and repairs canoes, the
canoe-maker is in reality devoting his labor to the taking of fish as
much as the actual fishermen, and the fish which he eats at night when
the fishermen come home are as truly the product of his labor as of
theirs. And thus when the division of labor is fairly inaugurated, and
instead of each attempting to satisfy all of his wants by direct resort
to nature, one fishes, another hunts, a third picks berries, a fourth
gathers fruit, a fifth makes tools, a sixth builds huts, and a seventh
prepares clothing---each one is to the extent he exchanges the direct
product of his own labor for the direct product of the labor of others
really applying his own labor to the production of the things he uses
---is in effect satisfying his particular desires by the exertion of his
particular powers; that is to say, what he receives he in reality
produces. If he digs roots and exchanges them for venison, he is in
effect as truly the procurer of the venison as though he had gone in
chase of the deer and left the huntsman to dig his own roots. The common
expression, ``I made so and so,'' signifying ``I earned so and so,'' or
``I earned money with which I purchased so and so,'' is, economically
speaking, not metaphorically but literally true. Earning is making.

Now, if we follow these principles, obvious enough in a simpler state of
society, through the complexities of the state we call civilized, we
shall see clearly that in every case in which labor is exchanged for
commodities, production really precedes enjoyment; that wages are the
earnings---that is to say, the makings of labor---not the advances of
capital, and that the laborer who receives his wages in money (coined or
printed, it may be, before his labor commenced) really receives in
return for the addition his labor has made to the general stock of
wealth, a draft upon that general stock, which he may utilize in any
particular form of wealth that will best satisfy his desires; and that
neither the money, which is but the draft, nor the particular form of
wealth which he uses it to call for, represents advances of capital for
his maintenance, but on the contrary represents the wealth, or a portion
of the wealth, his labor has already added to the general stock.

Keeping these principles in view we see that the draughtsman, who, shut
up in some dingy office on the banks of the Thames, is drawing the plans
for a great marine engine, is in reality devoting his labor to the
production of bread and meat as truly as though he were garnering the
grain in California or swinging a lariat on a La Plata pampa; that he is
as truly making his own clothing as though he were shearing sheep in
Australia or weaving cloth in Paisley, and just as effectually producing
the claret he drinks at dinner as though he gathered the grapes on the
banks of the Garonne. The miner who, two thousand feet under ground in
the heart of the Comstock, is digging out silver ore, is, in effect, by
virtue of a thousand exchanges, harvesting crops in valleys five
thousand feet nearer the earth's center; chasing the whale through
Arctic icefields; plucking tobacco leaves in Virginia; picking coffee
berries in Honduras; cutting sugar cane on the Hawaiian Islands;
gathering cotton in Georgia or weaving it in Manchester or Lowell;
making quaint wooden toys for his children in the Hartz Mountains; or
plucking amid the green and gold of Los Angeles orchards the oranges
which, when his shift is relieved, he will take home to his sick wife.
The wages which he receives on Saturday night at the mouth of the shaft,
what are they but the certificate to all the world that he has done
these things---the primary exchange in the long series which transmutes
his labor into the things he has really been laboring for?

All this is clear when looked at in this way; but to meet this fallacy
in all its strongholds and lurking places we must change our
investigation from the deductive to the inductive form. Let us now see,
if, beginning with facts and tracing their relations, we arrive at the
same conclusions as are thus obvious when, beginning with first
principles, we trace their exemplification in complex facts.

\newpage

\hypertarget{Book1Chapter2}{%
\section{CHAPTER II}\label{Book1Chapter2}}

\hypertarget{the-meaning-of-the-terms}{%
\subsection{THE MEANING OF THE TERMS}\label{the-meaning-of-the-terms}}

Before proceeding further in our inquiry, let us make sure of the
meaning of our terms, for indistinctness in their use must inevitably
produce ambiguity and indeterminateness in reasoning. Not only is it
requisite in economic reasoning to give to such words as ``wealth,''
``capital,'' ``rent,'' ``wages,'' and the like, a much more definite
sense than they bear in common discourse, but, unfortunately, even in
political economy there is, as to some of these terms, no certain
meaning assigned by common consent, different writers giving to the same
term different meanings, and the same writers often using a term in
different senses. Nothing can add to the force of what has been said by
so many eminent authors as to the importance of clear and precise
definitions, save the example, not an infrequent one, of the same
authors falling into grave errors from the very cause they warned
against. And nothing so shows the importance of language in thought as
the spectacle of even acute thinkers basing important conclusions upon
the use of the same word in varying senses. I shall endeavor to avoid
these dangers. It will be my effort throughout, as any term becomes of
importance, to state clearly what I mean by it, and to use it in that
sense and in no other. Let me ask the reader to note and to bear in mind
the definitions thus given, as otherwise I cannot hope to make myself
properly understood. I shall not attempt to attach arbitrary meanings to
words, or to coin terms, even when it would be convenient to do so, but
shall conform to usage as closely as is possible, only endeavoring so to
fix the meaning of words that they may clearly express thought.

What we have now on hand is to discover whether, as a matter of fact,
wages are drawn from capital. As a preliminary, let us settle what we
mean by wages and what we mean by capital. To the former word a
sufficiently definite meaning has been given by economic writers, but
the ambiguities which have attached to the use of the latter in
political economy will require a detailed examination.

As used in common discourse ``wages'' means a compensation paid to a
hired person for his services; and we speak of one man ``working for
wages,'' in contradistinction to another who is ``working for himself.''
The use of the term is still further narrowed by the habit of applying
it solely to compensation paid for manual labor. We do not speak of the
wages of professional men, managers or clerks, but of their fees,
commissions, or salaries. Thus the common meaning of the word wages is
the compensation paid to a hired person for manual labor. But in
political economy the word wages has a much wider meaning, and includes
all returns for exertion. For, as political economists explain, the
three agents or factors in production are land, labor, and capital, and
that part of the produce which goes to the second of these factors is by
them styled wages.

Thus the term labor includes all human exertion in the production of
wealth, and wages, being that part of the produce which goes to labor,
includes all reward for such exertion. There is, therefore, in the
politico-economic sense of the term wages no distinction as to the kind
of labor, or as to whether its reward is received through an employer or
not, but wages means the return received for the exertion of labor, as
distinguished from the return received for the use of capital, and the
return received by the landholder for the use of land. The man who
cultivates the soil for himself receives his wages in its produce, just
as, if he uses his own capital and owns his own land, he may also
receive interest and rent; the hunter's wages are the game he kills; the
fisherman's wages are the fish he takes. The gold washed out by the
self- employing gold-digger is as much his wages as the money paid to
the hired coal miner by the purchaser of his labor, and, as Adam Smith
shows, the high profits of retail storekeepers are in large part wages,
being the recompense of their labor and not of their capital. In short,
whatever is received as the result or reward of exertion is ``wages.''

This is all it is now necessary to note as to ``wages,'' but it is
important to keep this in mind. For in the standard economic works this
sense of the term wages is recognized with greater or less clearness
only to be subsequently ignored.

But it is more difficult to clear away from the idea of capital the
ambiguities that beset it, and to fix the scientific use of the term. In
general discourse, all sorts of things that have a value or will yield a
return are vaguely spoken of as capital, while economic writers vary so
widely that the term can hardly be said to have a fixed meaning. Let us
compare with each other the definitions of a few representative writers:

``That part of a man's stock,'' says Adam Smith (Book II, Chap. I),
``which he expects to afford him a revenue, is called his capital,'' and
the capital of a country or society, he goes on to say, consists of (1)
machines and instruments of trade which facilitate and abridge labor;
(2) buildings, not mere dwellings, but which may be considered
instruments of trade-such as shops, farmhouses, etc.; (3) improvements
of land which better fit it for tillage or culture; (4) the acquired and
useful abilities of all the inhabitants; (5) money; (6) provisions in
the hands of producers and dealers, from the sale of which they expect
to derive a profit; (7) the material of, or partially completed,
manufactured articles still in the hands of producers or dealers; (8)
completed articles still in the hands of producers or dealers. The first
four of these he styles fixed capital, and the last four circulating
capital, a distinction of which it is not necessary to our purpose to
take any note.

Ricardo's definition is:

\begin{quote}
Capital is that part of the wealth of a country which is employed in
production, and consists of food, clothing, tools, raw materials,
machinery, etc., necessary to give effect to labor.

\begin{itemize}
\tightlist
\item
  Principles of Political Economy, Chapter V.
\end{itemize}
\end{quote}

This definition, it will be seen, is very different from that of Adam
Smith, as it excludes many of the things which he includes---as acquired
talents, articles of mere taste or luxury in the possession of producers
or dealers; and includes some things he excludes such as food, clothing,
etc., in the possession of the consumer.

McCulloch's definition is:

\begin{quote}
The capital of a nation really comprises all those portions of the
produce of industry existing in it that may be directly employed either
to support human existence or to facilitate production.

\begin{itemize}
\tightlist
\item
  Notes on Wealth of Nations, Book II, Chap. I.
\end{itemize}
\end{quote}

This definition follows the line of Ricardo's, but is wider. While it
excludes everything that is not capable of aiding production, it
includes everything that is so capable, without reference to actual use
or necessity for use the horse drawing a pleasure carriage being,
according to McCulloch's view, as he expressly states, as much capital
as the horse drawing a plow, because he may, if need arises, be used to
draw a plow.

John Stuart Mill, following the same general line as Ricardo and
McCulloch, makes neither the use nor the capability of use, but the
determination to use, the test of capital. He says:

\begin{quote}
Whatever things are destined to supply productive labor with the
shelter, protection, tools and materials which the work requires, and to
feed and otherwise maintain the laborer during the process, are capital.

\begin{itemize}
\tightlist
\item
  Principles of Political Economy, Book I, Chap. IV.
\end{itemize}
\end{quote}

These quotations sufficiently illustrate the divergence of the masters.
Among minor authors the variance is still greater, as a few examples
will suffice to show.

Professor Wayland, whose ``Elements of Political Economy'' has long been
a favorite text-book in American educational institutions, where there
has been any pretense of teaching political economy, gives this lucid
definition:

\begin{quote}
The word capital is used in two senses. In relation to product it means
any substance on which industry is to be exerted. In relation to
industry, the material on which industry is about to confer value, that
on which it has conferred value; the instruments which are used for the
conferring of value, as well as the means of sustenance by which the
being is supported while he is engaged in performing the operation.''

\begin{itemize}
\tightlist
\item
  Elements of Political Economy, Book I, Chap. I.
\end{itemize}
\end{quote}

Henry C. Carey, the American apostle of protectionism, defines capital
as ``the instrument by which man obtains mastery over nature, including
in it the physical and mental powers of man himself.'' Professor Perry,
a Massachusetts free trader, very properly objects to this that it
hopelessly confuses the boundaries between capital and labor, and then
himself hopelessly confuses the boundaries between capital and land by
defining capital as ``any valuable thing outside of man himself from
whose use springs a pecuniary increase or profit.'' An English economic
writer of high standing, Mr.~Wm. Thornton, begins an elaborate
examination of the relations of labor and capital (``On Labor'') by
stating that he will include land with capital, which is very much as if
one who proposed to teach algebra should begin with the declaration that
he would consider the signs plus and minus as meaning the same thing and
having the same value. An American writer, also of high standing,
Professor Francis A. Walker, makes the same declaration in his elaborate
book on ``The Wages Question.'' Another English writer, N. A. Nicholson
(``The Science of Exchanges,'' London, 1873), seems to cap the climax of
absurdity by declaring in one paragraph (p. 26) that ``capital must of
course be accumulated by saving,'' and in the very next paragraph
stating that ``the land which produces a crop, the plow which turns the
soil, the labor which secures the produce, and the produce itself, if a
material profit is to be derived from its employment, are all alike
capital.'' But how land and labor are to be accumulated by saving them
he nowhere condescends to explain. In the same way a standard American
writer, Professor Amasa Walker (p.~66, ``Science of Wealth''), first
declares that capital arises from the net savings of labor and then
immediately afterward declares that land is capital.

I might go on for pages, citing contradictory and self-contradictory
definitions. But it would only weary the reader. It is unnecessary to
multiply quotations. Those already given are sufficient to show how wide
a difference exists as to the comprehension of the term capital. Any one
who wants further illustration of the ``confusion worse confounded''
which exists on this subject among the professors of political economy
may find it in any library where the works of these professors are
ranged side by side.

Now, it makes little difference what name we give to things, if when we
use the name we always keep in view the same things and no others. But
the difficulty arising in economic reasoning from these vague and
varying definitions of capital is that it is only in the premises of
reasoning that the term is used in the peculiar sense assigned by the
definition, while in the practical conclusions that are reached it is
always used, or at least it is always understood, in one general and
definite sense. When, for instance, it is said that wages are drawn from
capital, the word capital is understood in the same sense as when we
speak of the scarcity or abundance, the increase or decrease, the
destruction or increment, of capital-a commonly understood and definite
sense which separates capital from the other factors of production, land
and labor, and also separates it from like things used merely for
gratification. In fact, most people understand well enough what capital
is until they begin to define it, and I think their works will show that
the economic writers who differ so widely in their definitions use the
term in this commonly understood sense in all cases except in their
definitions and the reasoning based on them.

This common sense of the term is that of wealth devoted to procuring
more wealth. Dr.~Adam Smith correctly expresses this common idea when he
says: ``That part of a man's stock which he expects to afford him
revenue is called his capital.'' And the capital of a community is
evidently the sum of such individual stocks, or that part of the
aggregate stock which is expected to procure more wealth. This also is
the derivative sense of the term. The word capital, as philologists
trace it, comes down to us from a time when wealth was estimated in
cattle, and a man's income depended upon the number of head he could
keep for their increase.

The difficulties which beset the use of the word capital, as an exact
term, and which are even more strikingly exemplified in current
political and social discussions than in the definitions of economic
writers, arise from two facts-first, that certain classes of things, the
possession of which to the individual is precisely equivalent to the
possession of capital, are not part of the capital of the community;
and, second, that things of the same kind may or may not be capital,
according to the purpose to which they are devoted.

With a little care as to these points, there should be no difficulty in
obtaining a sufficiently clear and fixed idea of what the term capital
as generally used properly includes; such an idea as will enable us to
say what things are capital and what are not, and to use the word
without ambiguity or slip.

Land, labor and capital are the three factors of production. If we
remember that capital is thus a term used in contradistinction to the
land and labor, we at once see that nothing property included under
either one of these terms can be properly classed as capital. The term
land necessarily includes, not merely the surface of the earth as
distinguished from the water and air, but the whole material universe
outside of man himself, for it is only by having access to land, from
which his very body is drawn that man can come in contact with or use
nature. The term land embraces, in short, all natural materials, forces,
and opportunities, and, therefore, nothing that is freely supplied by
nature can be properly class as capital. A fertile field, a rich vein or
ore, a falling stream which supplies power, may give to the possessor
advantages equivalent to the possession of capital but to class such
things as capital would be to put an end to the distinction between land
and capital, and so far as they relate to each other, to make the two
terms meaningless. The term labor, in like manner, includes all human
exertion, and hence human powers whether natural or acquired can never
properly be classed as capital. In common parlance we often speak of a
man's knowledge, skill, or industry as constituting his capital; but
this is evidently a metaphorical use of language that must be eschewed
in reasoning that aims at exactness. Superiority in such qualities may
augment the income of an individuals just as capital would, and an
increase in knowledge, skill, or industry of a community may have the
same effect in increasing its production as would an increase in
capital; but this effect is due to the increased power of labor and not
to capital. Increased velocity may give to the impact of a cannon ball
the same effect as increased weight, yet, nevertheless, weight is one
thing and velocity another.

Thus, we must exclude from the category of capital everything that may
be included either as land or labor. Doing so, there remain only things
which are neither land nor labor, but which have resulted from the union
of these two original factors of production. Nothing can properly be
capital that does not consist of these - that is to say, nothing can be
capital that is not wealth.

But it is from ambiguities in the use of this inclusive term wealth that
many of the ambiguities which beset the term capital are derived.

As commonly used the word ``wealth'' is applied to anything having an
exchange value. But when used as a term of political economy it must be
limited to a much more definite meaning, because many things are
commonly spoken of as wealth which in taking account of collective or
general wealth cannot be considered as wealth at all. Such things have
an exchange value, and are commonly spoken of as wealth, insomuch as
they represent as between individuals, or between sets of individuals,
the power of obtaining wealth; but they are not truly wealth, inasmuch
as their increase or decrease does not affect the sum of wealth. Such as
bonds, mortgages, promissory notes, bank bills, or other stipulations
for the transfer of wealth. Such are slave, whose value represents
merely the power of one class to appropriate the earnings of another
class. Such are lands, or other natural opportunities, the value of
which is but the result of the acknowledgment in favor of certain
persons of an exclusive right to their use, and which represents merely
the power thus given to the owners to demand a share of the wealth
produced by those who use them. Increase in the amount of bonds,
mortgages, notes or bank bills, cannot increase the wealth of the
community that includes as well those who promise to pay as those who
are entitled to receive. The enslavement of a part of their number could
not increase the wealth of a people, for what the enslavers gained the
enslaved would lose. Increase in land values does not represent increase
in the common wealth, for what land owners gain by higher prices, the
tenants or purchasers who must pay them will lose. And all this relative
wealth, which in common thought and speech, in legislation and law, is
undistinguished from actual wealth, could, without the destruction or
consumption for anything more than a few drops of ink and a piece of
paper, be utterly annihilated. By enactment of the sovereign political
power debts might be canceled, slaves emancipated, and land resumed as
the common property of the whole people, without the aggregate wealth
being diminished by the value of a pinch of snuff, for what some would
lose others would gain. There would be no more destruction of wealth
than there was creation of wealth when Elizabeth Tudor enriched her
favorite courtiers by the grant of monopolies, or when Boris Godunof
made Russian peasants merchantable property.

All things which have an exchange value are, therefore, not wealth, in
the only sense in which the term can be used in political economy. Only
such things can be wealth the production of which increase and the
destruction of which decreases the aggregate of wealth. If we consider
what these things are, and what their nature is, we shall have no
difficulty in defining wealth.

When we speak of a community increasing in wealth---as when we say that
England has increased in wealth since the accession of Victoria, or that
California is a wealthier country than when it was a Mexican
territory-we do not mean to say that there is more land, or that the
natural powers of the land are greater, or that there are more people,
for when we wish to express that idea we speak of increase of
population; or that the debts or dues owing by some of these people to
others of their number have increased; but we mean that there is an
increase of certain tangible things, having an actual and not merely a
relative value-such as buildings, cattle, tools, machinery, agricultural
and mineral products, manufactured goods, ships, wagons, furniture, and
the like. The increase of such things constitutes an increase of wealth;
their decrease is a lessening of wealth; and the community that, in
proportion to its numbers, has most of such things is the wealthiest
community. The common character of these things is that they consist of
natural substances or products which have been adapted by human labor to
human use or gratification, their value depending on the amount of labor
which upon the average would be required to produce things of like kind.

Thus wealth, as alone the term can be used in political economy,
consists of natural products that have been secured, moved, combined,
separated, or in other ways modified by human exertion, so as to fit
them for the gratification of human desires. It is, in other words,
labor impressed upon matter in such a way as to store up, as the heat of
the sun is stored up in coal, the power of human labor to minister to
human desires. Wealth is not the sole object of labor, for labor is also
expended in ministering directly to desire; but it is the object and
result of what we call productive labor---that is, labor which gives
value to material things. Nothing which nature supplies to man without
his labor is wealth, nor yet does the expenditure of labor result in
wealth unless there is a tangible product which has and retains the
power of ministering to desire.

Now, as capital is wealth devoted to a certain purpose, nothing can be
capital which does not fall within this definition of wealth. By
recognizing and keeping this in mind, we get rid of misconceptions which
vitiate all reasoning in which they are permitted, which befog popular
thought, and have led into mazes of contradiction even acute thinkers.

But though all capital is wealth, all wealth is not capital. Capital is
only a part of wealth-that part, namely, which is devoted to the aid of
production. It is in drawing this line between the wealth that is and
the wealth that is not capital that a second class of misconceptions are
likely to occur.

The errors which I have been pointing out, and which consist in
confounding with wealth and capital things essentially distinct, or
which have but a relative existence, are now merely vulgar errors. They
are widespread, it is true, and have a deep root, being held, not merely
by the less educated classes, but seemingly by a large majority of those
who in such advanced countries as England and the United States mold and
guide public opinion, make the laws in Parliaments, Congresses and
Legislatures, and administer them in the courts. They crop out,
moreover, in the disquisitions of many of those flabby writers who have
burdened the press and darkened counsel by numerous volumes which are
dubbed political economy, and which pass as text-books with the ignorant
and as authority with those who do not think for themselves.
Nevertheless, they are only vulgar errors, inasmuch as they receive no
countenance from the best writers on political economy. By one of those
lapses which flaw his great work and strikingly evince the imperfections
of the highest talent, Adam Smith counts as capital certain personal
qualities, an inclusion which is not consistent with his original
definition of capital as stock from which revenue is expected. But this
error has been avoided by his most eminent successors, and in the
definitions, previously given, of Ricardo, McCulloch, and Mill, it is
not involved. Neither in their definitions nor in that of Smith is
involved the vulgar error which confounds as real capital things which
are only relatively capital, such as evidences of debt, land values,
etc. But as to things which are really wealth, their definitions differ
from each other, and widely from that of Smith, as to what is and what
is not to be considered as capital. The stock of a jeweler would, for
instance, be included as capital by the definition of Smith, and the
food or clothing in possession of a laborer would be excluded. But the
definitions of Ricardo and McCulloch would exclude the stock of the
jeweler, as would also that of Mill, if understood as most persons would
understand the words I have quoted. But as explained by him, it is
neither the nature nor the destination of the things themselves which
determines whether they are or are not capital, but the intention of the
owner to devote either the things or the value received from their sale
to the supply of productive labor with tools, materials, and
maintenance. All these definitions, however, agree in including as
capital the provisions and clothing of the laborer, which Smith
excludes.

Let us consider these three definitions, which represent the best
teachings of current political economy:

To McCulloch's definition of capital as ``all those portions of the
produce of industry that may be directly employed either to support
human existence or to facilitate production,'' there are obvious
objections. One may pass along any principal street in a thriving town
or city and see stores filled with all sorts of valuable things, which,
though they cannot be employed either to support human existence or to
facilitate production, undoubtedly constitute part of the capital of the
storekeepers and part of the capital of the community. And he can also
see products of industry capable of supporting human existence or
facilitating production being consumed in ostentation or useless luxury.
Surely these, though they might, do not constitute part of capital.

Ricardo's definition avoids including as capital things which might be
but are not employed in production, by covering only such as are
employed. But it is open to the first objection made to McCulloch's. If
only wealth that may be, or that is, or that is destined to be, used in
supporting producers, or assisting production, is capital, then the
stocks of jewelers, toy dealers, tobacconists, confectioners, picture
dealers, etc.---in fact, all stocks that consist of, and all stocks in
so far as they consist of articles of luxury, are not capital.

If Mill, by remitting the distinction to the mind of the capitalist,
avoids this difficulty (which does not seem to me clear), it is by
making the distinction so vague that no power short of omniscience could
tell in any given country at any given time what was and what was not
capital.

But the great defect which these definitions have in common is that they
include what clearly cannot be accounted capital, if any distinction is
to be made between laborer and capitalist. For they bring into the
category of capital the food, clothing, etc., in the possession of the
day laborer, which he will consume whether he works or not, as well as
the stock in the hands of the capitalist, with which he proposes to pay
the laborer for his work.

Yet, manifestly, this is not the sense in which the term capital is used
by these writers when they speak of labor and capital as taking separate
parts in the work of production and separate shares in the distribution
of its proceeds; when they speak of wages as drawn from capital, or as
depending upon the ratio between labor and capital, or in any of the
ways in which the term is generally used by them. In all these cases the
term capital is used in its commonly understood sense, as that portion
of wealth which its owners do not propose to use directly for their own
gratification, but for the purpose of obtaining more wealth. In short,
by political economists, in everything except their definitions and
first principles, as well as by the world at large, ``that part of a
man's stock,'' to use the words of Adam Smith, ``which he expects to
afford him revenue is called his capital.'' This is the only sense in
which the term capital expresses any fixed idea-the only sense in which
we can with any clearness separate it from wealth and contrast it with
labor. For, if we must consider as capital everything which supplies the
laborer with food, clothing, shelter, etc., then to find a laborer who
is not a capitalist we shall be forced to hunt up an absolutely naked
man, destitute even of a sharpened stick, or of a burrow in the
ground---a situation in which, save as the result of exceptional
circumstances, human beings have never yet been found.

It seems to me that the variance and inexactitude in these definitions
arise from the fact that the idea of what capital is has been deduced
from a preconceived idea of how capital assists production. Instead of
determining what capital is, and then observing what capital does, the
functions of capital have first been assumed, and then a definition of
capital made which includes all things which do or may perform those
functions. Let us reverse this process, and, adopting the natural order,
ascertain what the thing is before settling what it does. All we are
trying to do, all that it is necessary to do, is to fix, as it were, the
metes and bounds of a term that in the main is well apprehended to make
definite, that is, sharp and clear on its verges, a common idea.

If the articles of actual wealth existing at a given time in a given
community were presented in situ to a dozen intelligent men who had
never read a line of political economy, it is doubtful if they would
differ in respect to a single item, as to whether it should be accounted
capital or not. Money which its owner holds for use in his business or
in speculation would be accounted capital; money set aside for household
or personal expenses would not. That part of a farmer's crop held for
sale or for seed, or to feed his help in part payment of wages, would be
accounted capital; that held for the use of his own family would not be.
The horses and carriage of a hackman would be classed as capital, but an
equipage kept for the pleasure of its owner would not. So no one would
think of counting as capital the false hair on the head of a woman, the
cigar in the mouth of a smoker, or the toy with which a child is
playing; but the stock of a hair dealer, of a tobacconist, or of the
keeper of a toy store, would be unhesitatingly set down as capital. A
coat which a tailor had made for sale would be accounted capital, but
not the coat he had made for himself. Food in the possession of a
hotel-keeper or a restaurateur would be accounted capital, but not the
food in the pantry of a housewife, or in the lunch basket of a workman.
Pig iron in the hands of the smelter, or founder, or dealer, would be
accounted capital, but not the pig iron used as ballast in the hold of a
yacht. The bellows of a blacksmith, the looms of a factory, would be
capital, but not the sewing machine of a woman who does only her own
work; a building let for hire, or used for business or productive
purposes, but not a homestead. In short, I think we should find that
now, as when Dr.~Adam Smith wrote, ``that part of a man's stock which he
expects to yield him a revenue is called his capital.'' And, omitting
his unfortunate slip as to personal qualities, and qualifying somewhat
his enumeration of money, it is doubtful if we could better list the
different articles of capital than did Adam Smith in the passage which
in the previous part of this chapter I have condensed.

Now, if, after having thus separated the wealth that is capital from the
wealth that is not capital, we look for the distinction between the two
classes, we shall not find it to be as to the character, capabilities,
or final destination of the things themselves, as has been vainly
attempted to draw it; but it seems to me that we shall find it to be as
to whether they are or are not in the possession of the consumer. Such
articles of wealth as in themselves, in their uses, or in their
products, are yet to be exchanged are capital; such articles of wealth
as are in the hands of the consumer are not capital. Hence, if we define
capital as wealth in course of exchange, understanding exchange to
include not merely the passing from hand to hand, but also such
transmutations as occur when the reproductive or transforming forces of
nature are utilized for the increase of wealth, we shall, I think,
comprehend all the things that the general idea of capital properly
includes, and shut out all it does not. Under this definition, it seems
to me, for instance, will fall all such tools as are really capital. For
it is as to whether its services or uses are to be exchanged or not
which makes a tool an article of capital or merely an article of wealth.
Thus, the lathe of a manufacturer used in making things which are to be
exchanged is capital, while the lathe kept by a gentleman for his own
amusement is not. Thus, wealth used in the construction of a railroad, a
public telegraph line, a stage coach, a theater, a hotel, etc., may be
said to be placed in the course of exchange. The exchange is not
effected all at once, but little by little, with an indefinite number of
people. Yet there is an exchange, and the ``consumers'' of the railroad,
the telegraph line, the stage coach, theater or hotel, are not the
owners, but the persons who from time to time use them.

Nor is this definition inconsistent with the idea that capital is that
part of wealth devoted to production. It is too narrow an understanding
of production which confines it merely to the making of things.
Production includes not merely the making of things, but the bringing of
them to the consumer. The merchant or storekeeper is thus as truly a
producer as is the manufacturer, or farmer, and his stock or capital is
as much devoted to production as is theirs. But it is not worth while
now to dwell upon the functions of capital, which we shall be better
able to determine hereafter. Nor is the definition of capital I have
suggested of any importance. I am not writing a text-book, but only
attempting to discover the laws which control a great social problem,
and if the reader has been led to form a clear idea of what things are
meant when we speak of capital my purpose is served.

But before closing this digression let me call attention to what is
often forgotten---namely, that the terms ``wealth,'' ``capital,''
``wages,'' and the like, as used in political economy are abstract
terms, and that nothing can be generally affirmed or denied of them that
cannot be affirmed or denied of the whole class of things they
represent. The failure to bear this in mind has led to much confusion of
thought, and permits fallacies, otherwise transparent, to pass for
obvious truths. Wealth being an abstract term, the idea of wealth, it
must be remembered, involves the idea of exchange ability. The
possession of wealth to a certain amount is potentially the possession
of any or all species of wealth to that equivalent in exchange. And,
consequently, so of capital.

\newpage

\hypertarget{Book1Chapter3}{%
\section{CHAPTER III}\label{Book1Chapter3}}

\hypertarget{wages-not-drawn-from-capital-but-produced-by-the-labor}{%
\subsection{WAGES NOT DRAWN FROM CAPITAL, BUT PRODUCED BY THE
LABOR}\label{wages-not-drawn-from-capital-but-produced-by-the-labor}}

The importance of this digression will, I think, become more and more
apparent as we proceed in our inquiry, but its pertinency to the branch
we are now engaged in may at once be seen.

It is at first glance evident that the economic meaning of the term
wages is lost sight of, and attention is concentrated upon the common
and narrow meaning of the word, when it is affirmed that wages are drawn
from capital. For, in all those cases in which the laborer is his own
employer and takes directly the produce of his labor as its reward, it
is plain enough that wages are not drawn from capital, but result
directly as the product of the labor. If, for instance, I devote my
labor to gathering birds' eggs or picking wild berries, the eggs or
berries I thus get are my wages. Surely no one will contend that in such
a case wages are drawn from capital. There is no capital in the case. An
absolutely naked man, thrown on an island where no human being has
before trod, may gather birds' eggs or pick berries.

Or if I take a piece of leather and work it up into a pair of shoes, the
shoes are my wages---the reward of my exertion. Surely they are not
drawn from capital---either my capital or anyone else's capital---but
are brought into existence by the labor of which they become the wages;
and in obtaining this pair of shoes as the wages of my labor, capital is
not even momentarily lessened one iota. For, if we call in the idea of
capital, my capital at the beginning consists of the piece of leather,
the thread, etc. As my labor goes on, value is steadily added, until,
when my labor results in the finished shoes, I have my capital plus the
difference in value between the material and the shoes. In obtaining
this additional value---my wages---how is capital at any time drawn
upon?

Adam Smith, who gave the direction to economic thought that has resulted
in the current elaborate theories of the relation between wages and
capital, recognized the fact that in such simple cases as I have
instanced, wages are the produce of labor, and thus begins his chapter
upon the wages of labor (Chapter VIII):

\begin{quote}
The produce of labor constitutes the natural recompense or wages of
labor. In that original state of things which precedes both the
appropriation of land and the accumulation of stock, the whole produce
of labor belongs to the laborer. He has neither landlord nor master to
share with him.
\end{quote}

Had the great Scotchman taken this as the initial point of his
reasoning, and continued to regard the produce of labor as the natural
wages of labor, and the landlord and master but as sharers, his
conclusions would have been very different, and political economy today
would not embrace such a mass of contradictions and absurdities; but
instead of following the truth obvious in the simple modes of production
as a clew through the perplexities of the more complicated forms, he
momentarily recognizes it, only immediately to abandon it, and stating
that ``in every part of Europe twenty workmen serve under a master for
one that is independent,'' he recommences the inquiry from a point of
view in which the master is considered as providing from his capital the
wages of his workmen.

It is evident that in thus placing the proportion of self-employing
workmen as but one in twenty, Adam Smith had in mind but the mechanic
arts, and that, including all laborers, the proportion who take their
earnings directly, without the intervention of an employer, must, even
in Europe a hundred years ago, have been much greater than this. For,
besides the independent laborers who in every community exist in
considerable numbers, the agriculture of large districts of Europe has,
since the time of the Roman Empire, been carried on by the metayer
system, under which the capitalist receives his return from the laborer
instead of the laborer from the capitalist. At any rate, in the United
States, where any general law of wages must apply as fully as in Europe,
and where in spite of the advance of manufactures a very large part of
the people are yet self-employing farmers, the proportion of laborers
who get their wages through an employer must be comparatively small.

But it is not necessary to discuss the ratio in which self-employing
laborers anywhere stand to hired laborers, nor is it necessary to
multiply illustrations of the truism that where the laborer takes
directly his wages they are the product of his labor, for as soon as it
is realized that the term wages includes all the earnings of labor, as
well when taken directly by the laborer in the results of his labor as
when received from an employer, it is evident that the assumption that
wages are drawn from capital, on which as a universal truth such a vast
superstructure is in standard politico-economic treatises so
unhesitatingly built, is at least in large part untrue, and the utmost
that can with any plausibility be affirmed is that some wages (i.e.,
wages received by the laborer from an employer) are drawn from capital.
This restriction of the major premise at once invalidates all the
deductions that are made from it; but without resting here, let us see
whether even in this restricted sense it accords with the facts. Let us
pick up the clew where Adam Smith dropped it, and advancing step by
step, see whether the relation of facts which is obvious in the simplest
forms of production does not run through the most complex.

Next in simplicity to ``that original state of things,'' of which many
examples may yet be found, where the whole produce of labor belongs to
the laborer, is the arrangement in which the laborer, though working for
another person, or with the capital of another person, receives his
wages in kind-that is to say, in the things his labor produces. In this
case it is as clear as in the case of the self-employing laborer that
the wages are really drawn from the product of the labor, and not at all
from capital. If I hire a man to gather eggs, to pick berries, or to
make shoes, paying him from the eggs, the berries, or the shoes that his
labor secures, there can be no question that the source of the wages is
the labor for which they are paid. Of this form of hiring is the
saer-and-daer stock tenancy, treated of with such perspicuity by Sir
Henry Maine in his ``Early History of Institutions,'' and which so
clearly involved the relation of employer and employed as to render the
acceptor of cattle the man or vassal of the capitalist who thus employed
him. It was on such terms as these that Jacob worked for Laban, and to
this day, even in civilized countries, it is not an infrequent mode of
employing labor. The farming of land on shares, which prevails to a
considerable extent in the Southern States of the Union and in
California, the metayer system of Europe, as well as the many cases in
which superintendents, salesmen, etc., are paid by a percentage of
profits, what are they but the employment of labor for wages which
consist of part of its produce?

The next step in the advance from simplicity to complexity is where the
wages, though estimated in kind, are paid in an equivalent of something
else. For instance, on American whaling ships the custom is not to pay
fixed wages, but a ``lay,'' or proportion of the catch, which varies
from a sixteenth to a twelfth to the captain down to a three-hundredth
to the cabin-boy. Thus, when a whaleship comes into New Bedford or San
Francisco after a successful cruise, she carries in her hold the wages
of her crew, as well as the profits of her owners, and an equivalent
which will reimburse them for all the stores used up during the voyage.
Can anything be clearer than that these wages---this oil and bone which
the crew of the whaler have taken---have not been drawn from capital,
but are really a part of the produce of their labor? Nor is this fact
changed or obscured in the slightest degree where, as a matter of
convenience, instead of dividing up between the crew their proportion of
the oil and bone, the value of each man's share is estimated at the
market price, and he is paid for it in money. The money is but the
equivalent of the real wages, the oil and bone. In no way is there any
advance of capital in this payment. The obligation to pay wages does not
accrue until the value from which they are to be paid is brought into
port. At the moment when the owner takes from his capital money to pay
the crew he adds to his capital oil and bone.

So far there can be no dispute. Let us now take another step, which will
bring us to the usual method of employing labor and paying wages.

The Farallone Islands, off the Bay of San Francisco, are a hatching
ground of sea-fowl, and a company who claim these islands employ men in
the proper season to collect the eggs. They might employ these men for a
proportion of the eggs they gather, as is done in the whale fishery, and
probably would do so if there were much uncertainty attending the
business; but as the fowl are plentiful and tame, and about so many eggs
can be gathered by so much labor, they find it more convenient to pay
their men fixed wages. The men go out and remain on the islands,
gathering the eggs and bringing them to a landing, whence, at intervals
of a few days, they are taken in a small vessel to San Francisco and
sold. When the season is over the men return and are paid their
stipulated wages in coin. Does not this transaction amount to the same
thing as if, instead of being paid in coin, the stipulated wages were
paid in an equivalent of the eggs gathered? Does not the coin represent
the eggs, by the sale of which it was obtained, and are not these wages
as much the product of the labor for which they are paid as the eggs
would be in the possession of a man who gathered them for himself
without the intervention of any employer?

To take another example, which shows by reversion the identity of wages
in money with wages in kind. In San Buenaventura lives a man who makes
an excellent living by shooting for their oil and skins the common hair
seals which frequent the islands forming the Santa Barbara Channel. When
on these sealing expeditions he takes two or three Chinamen along to
help him, whom at first he paid wholly in coin. But it seems that the
Chinese highly value some of the organs of the seal, which they dry and
pulverize for medicine, as well as the long hairs in the whiskers of the
male seal, which, when over a certain length, they greatly esteem for
some purpose that to outside barbarians is not very clear. And this man
soon found that the Chinamen were very willing to take instead of money
these parts of the seals killed, so that now, in large part, he thus
pays them their wages.

Now, is not what may be seen in all these cases---the identity of wages
in money with wages in kind---true of all cases in which wages are paid
for productive labor? Is not the fund created by the labor really the
fund from which the wages are paid?

It may, perhaps, be said: ``There is this difference where a man works
for himself, or where, when working for an employer, he takes his wages
in kind, his wages depend upon the result of his labor. Should that,
from any misadventure, prove futile, he gets nothing. When he works for
an employer, however, he gets his wages anyhow-they depend upon the
performance of the labor, not upon the result of the labor.'' But this
is evidently not a real distinction. For on the average, the labor that
is rendered for fixed wages not only yields the amount of the wages, but
more; else employers could make no profit. When wages are fixed, the
employer takes the whole risk and is compensated for this assurance, for
wages when fixed are always somewhat less than wages contingent. But
though when fixed wages are stipulated the laborer who has performed his
part of the contract has usually a legal claim upon the employer, it is
frequently, if not generally, the case that the disaster which prevents
the employer from reaping benefit from the labor prevents him from
paying the wages. And in one important department of industry the
employer is legally exempt in case of disaster, although the contract be
for wages certain and not contingent. For the maxim of admiralty law is,
that ``freight is the mother of wages,'' and though the seaman may have
performed his part, the disaster which prevents the ship from earning
freight deprives him of claim for his wages.

In this legal maxim is embodied the truth for which I am contending.
Production is always the mother of wages. Without production, wages
would not and could not be. It is from the produce of labor, not from
the advances of capital that wages come.

Wherever we analyze the facts this will be found to be true. For labor
always precedes wages. This is as universally true of wages received by
the laborer from an employer as it is of wages taken directly by the
laborer who is his own employer. In the one class of cases as in the
other, reward is conditioned upon exertion. Paid sometimes by the day,
oftener by the week or month, occasionally by the year, and in many
branches of production by the piece; the payment of wages by an employer
to an employee always implies the previous rendering of labor by the
employee for the benefit of the employer, for the few cases in which
advance payments are made for personal services are evidently referable
either to charity or to guarantee and purchase. The name ``retainer,''
given to advance payments to lawyers, shows the true character of the
transaction, as does the name ``blood money'' given in 'longshore
vernacular to a payment which is nominally wages advanced to sailors,
but which in reality is purchase money---both English and American law
considering a sailor as much a chattel as a pig.

I dwell on this obvious fact that labor always precedes wages, because
it is all-important to an understanding of the more complicated
phenomena of wages that it should be kept in mind. And obvious as it is,
as I have put it, the plausibility of the proposition that wages are
drawn from capital-a proposition that is made the basis for such
important and far-reaching deductions-comes in the first instance from a
statement that ignores and leads the attention away from this truth.
That statement is, that labor cannot exert its productive power unless
supplied by capital with maintenance.* The unwary reader at once
recognizes the fact that the laborer must have food, clothing, etc., in
order to enable him to perform the work, and having been told that the
food, clothing, etc., used by productive laborers are capital, he
assents to the conclusion that the consumption of capital is necessary
to the application of labor, and from this it is but an obvious
deduction that industry is limited by capital-that the demand for labor
depends upon the supply of capital, and hence that wages depend upon the
ratio between the number of laborers looking for employment and the
amount of capital devoted to hiring them.

\begin{itemize}
\item
  Industry is limited by capital. . . . There can be no more industry
  than is supplied with materials to work up and food to eat.
  Self-evident as the thing is, it is often forgotten that the people of
  a country are maintained and have their wants supplied not by the
  produce of present labor, but of past. They consume what has been
  produced, not what is about to be produced. Now, of what has been
  produced a part only is allotted to the support of productive labor,
  and there will not and cannot be more of that labor than the portion
  so allotted (which is the capital of the country) can feed and provide
  with the materials and instruments of production.

  \begin{itemize}
  \tightlist
  \item
    John Stuart Mill, Principles of Political Economy, Book I
  \end{itemize}
\end{itemize}

But I think the discussion in the previous chapter will enable anyone to
see wherein lies the fallacy of this reasoning-a fallacy which has
entangled some of the most acute minds in a web of their own spinning.
It is in the use of the term capital in two senses. In the primary
proposition that capital is necessary to the exertion of productive
labor, the term ``capital'' is understood as including all food,
clothing, shelter, etc.; whereas, in the deductions finally drawn from
it, the term is used in its common and legitimate meaning of wealth
devoted, not to the immediate gratification of desire, but to the
procurement of more wealth---of wealth in the hands of employers as
distinguished from laborers. The conclusion is no more valid than it
would be from the acceptance of the proposition that a laborer cannot go
to work without his breakfast and some clothes, to infer that no more
laborers can go to work than employers first furnish with breakfasts and
clothes. Now, the fact is that laborers generally furnish their own
breakfasts and the clothes in which they go to work; and the further
fact is, that capital (in the sense in which the word is used in
distinction to labor) in exceptional cases sometimes may, but is never
compelled to make advances to labor before the work begins. Of all the
vast number of unemployed laborers in the civilized world to-day, there
is probably not a single one willing to work who could not be employed
without any advance of wages. A great proportion would doubtless gladly
go to work on terms which did not require the payment of wages before
the end of a month; it is doubtful if there are enough to be called a
class who would not go to work and wait for their wages until the end of
the week, as most laborers habitually do; while there are certainly none
who would not wait for their wages until the end of the day, or if you
please, until the next meal hour. The precise time of the payment of
wages is immaterial; the essential point---the point I lay stress
on---is that it is after the performance of work.

The payment of wages, therefore, always implies the previous rendering
of labor. Now, what does the rendering of labor in production imply?
Evidently the production of wealth, which, if it is to be exchanged or
used in production, is capital. Therefore, the payment of capital in
wages pre-supposes a production of capital by the labor for which the
wages are paid. And as the employer generally makes a profit, the
payment of wages is, so far as he is concerned, but the return to the
laborer of a portion of the capital he has received from the labor. So
far as the employee is concerned, it is but the receipt of a portion of
the capital his labor has previously produced. As the value paid in the
wages is thus exchanged for a value brought into being by the labor, how
can it be said that wages are drawn from capital or advanced by capital?
As in the exchange of labor for wages the employer always gets the
capital created by the labor before he pays out capital in the wages, at
what point is his capital lessened even temporarily?*

\begin{itemize}
\tightlist
\item
  I speak of labor producing capital for the sake of greater clearness.
  What labor always procures is either wealth, which may or may not be
  capital, or services, the cases in which nothing is obtained being
  merely exceptional cases of misadventure. Where the object of the
  labor is simply the gratification of the employer, as where I hire a
  man to black my boots, I do not pay the wages from capital, but from
  wealth which I have devoted, not to reproductive uses, but to
  consumption for my own satisfaction. Even if wages thus paid be
  considered as drawn from capital, then by that act they pass from the
  category of capital to that of wealth devoted to the gratification of
  the possessor, as when a cigar dealer takes a dozen cigars from the
  stock he has for sale and puts them in his pocket for his own use.
\end{itemize}

Bring the question to the test of facts. Take, for instance, an
employing manufacturer who is engaged in turning raw material into
finished products-cotton into cloth, iron into hardware, leather into
boots, or so on, as may be, and who pays his hands, as is generally the
case, once a week. Make an exact inventory of his capital on Monday
morning before the beginning of work, and it will consist of his
buildings, machinery, raw materials, money on hand, and finished
products in stock. Suppose, for the sake of simplicity, that he neither
buys nor sells during the week, and after work has stopped and he has
paid his hands on Saturday night, take a new inventory of his capital.
The item of money will be less, for it has been paid out in wages; there
will be less raw material, less coal, etc., and a proper deduction must
be made from the value of the buildings and machinery for the week's
wear and tear. But if he is doing a remunerative business, which must on
the average be the case, the item of finished products will be so much
greater as to compensate for all these deficiencies and show in the
summing up an increase of capital. Manifestly, then, the value he paid
his hands in wages was not drawn from his capital, or from any one
else's capital. It came, not from capital, but from the value created by
the labor itself. There was no more advance of capital than if he had
hired his hands to dig clams, and paid them with a part of the clams
they dug. Their wages were as truly the produce of their labor as were
the wages of the primitive man, when, long ``before the appropriation of
land and the accumulation of stock,'' he obtained an oyster by knocking
it with a stone from the rocks.

As the laborer who works for an employer does not get his wages until he
has performed the work, his case is similar to that of the depositor in
a bank who cannot draw money out until he has put money in. And as by
drawing out what he has previously put in, the bank depositor does not
lessen the capital of the bank, neither can laborers by receiving wages
lessen even temporarily either the capital of the employer or the
aggregate capital of the community. Their wages no more come from
capital than the checks of depositors are drawn against bank capital. It
is true that laborers in receiving wages do not generally receive back
wealth in the same form in which they have rendered it, any more than
bank depositors receive back the identical coins or bank notes they have
deposited, but they receive it in equivalent form, and as we are
justified in saying that the depositor receives from the bank the money
he paid in, so are we justified in saying that the laborer receives in
wages the wealth he has rendered in labor.

That this universal truth is so often obscured, is largely due to that
fruitful source of economic obscurities, the confounding of wealth with
money; and it is remarkable to see so many of those who, since Dr.~Adam
Smith made the egg stand on its head, have copiously demonstrated the
fallacies of the mercantile system, fall into delusions of the very same
kind in treating of the relations of capital and labor. Money being the
general medium of exchanges, the common flux through which all
transmutations of wealth from one form to another take place, whatever
difficulties may exist to an exchange will generally show themselves on
the side of reduction to money, and thus it is sometimes easier to
exchange money for any other form of wealth than it is to exchange
wealth in a particular form into money, for the reason that there are
more holders of wealth who desire to make some exchange than there are
who desire to make any particular exchange. And so a producing employer
who has paid out his money in wages may sometimes find it difficult to
turn quickly back into money the increased value for which his money has
really been exchanged, and is spoken of as having exhausted or advanced
his capital in the payment of wages. Yet, unless the new value created
by the labor is less than the wages paid, which can be only an
exceptional case, the capital which he had before in money he now has in
goods-it has been changed in form, but not lessened.

There is one branch of production in regard to which the confusions of
thought which arise from the habit of estimating capital in money are
least likely to occur, inasmuch as its product is the general material
and standard of money. And it so happens that this business furnishes
us, almost side by side, with illustrations of production passing from
the simplest to most complex forms.

In the early days of California, as afterward in Australia, the placer
miner, who found in river bed or surface deposit the glittering
particles which the slow processes of nature had for ages been
accumulating, picked up or washed out his ``wages'' (so, too, he called
them) in actual money, for coin being scarce, gold dust passed as
currency by weight, and at the end of the day had his wages in money in
a buckskin bag in his pocket. There can be no dispute as to whether
these wages came from capital or not. They were manifestly the produce
of his labor. Nor could there be any dispute when the holder of a
specially rich claim hired men to work for him and paid them off in the
identical money which their labor had taken from gulch or bar. As coin
became more abundant, its greater convenience in saving the trouble and
loss of weighing assigned gold dust to the place of a commodity, and
with coin obtained by the sale of the dust their labor had procured, the
employing miner paid off his hands. Where he had coin enough to do so,
instead of selling his gold dust at the nearest store and paying a
dealer's profit, he retained it until he got enough to take a trip, or
send by express to San Francisco, where at the mint he could have it
turned into coin without charge. While thus accumulating gold dust he
was lessening his stock of coin; just as the manufacturer, while
accumulating a stock of goods, lessens his stock of money. Yet no one
would be obtuse enough to imagine that in thus taking in gold dust and
paying out coin the miner was lessening his capital.

But the deposits that could be worked without preliminary labor were
soon exhausted, and gold mining rapidly took a more elaborate character.
Before claims could be opened so as to yield any return deep shafts had
to be sunk, great dams constructed, long tunnels out through the hardest
rock, water brought for miles over mountain ridges and across·deep
valleys, and expensive machinery put up. These works could not be
constructed without capital. Sometimes their construction required
years, during which no return could be hoped for, while the men employed
had to be paid their wages every week, or every month. Surely, it will
be said, in such cases, even if in no others, that wages do actually
come from capital; are actually advanced by capital; and must
necessarily lessen capital in their payment! Surely here, at least,
industry is limited by capital, for without capital such works could not
be carried on! Let us see:

It is cases of this class that are always instanced as showing that
wages are advanced from capital. For where wages are paid before the
object of the labor is obtained, or is finished-as in agriculture, where
plowing and sowing must precede by several months the harvesting of the
crop; as in the erection of buildings, the construction of ships,
railroads, canals, etc.---it is clear that the owners of the capital
paid in wages cannot expect an immediate return, but, as the phrase is,
must ``outlay it,'' or ``lie out of it'' for a time, which sometimes
amounts to many years. And hence, if first principles are not kept in
mind, it is easy to jump to the conclusion that wages are advanced by
capital.

But such cases will not embarrass the reader to whom in what has
preceded I have made myself clearly understood. An easy analysis will
show that these instances where wages are paid before the product is
finished, or even produced, do not afford any exception to the rule
apparent where the product is finished before wages are paid.

If I go to a broker to exchange silver for gold, I lay down my silver,
which he counts and puts away, and then hands me the equivalent in gold,
minus his commission. Does the broker advance me any capital? Manifestly
not. What he had before in gold he now has in silver, plus his profit.
And as he got the silver before he paid out the gold, there is on his
part not even momentarily an advance of capital.

Now, this operation of the broker is precisely analogous to what the
capitalist does, when, in such cases as we are now considering, he pays
out capital in wages. As the rendering of labor precedes the payment of
wages, and as the rendering of labor in production implies the creation
of value, the employer receives value before he pays out value---he but
exchanges capital of one form for capital of another form. For the
creation of value does not depend upon the finishing of the product; it
takes place at every stage of the process of production, as the
immediate result of the application of labor, and hence, no matter how
long the process in which it is engaged, labor always adds to capital by
its exertion before it takes from capital in its wages.

Here is a blacksmith at his forge making picks. Clearly he is making
capital-adding picks to his employer's capital before he draws money
from it in wages. Here is a machinist or boilermaker working on the
keel-plates of a Great Eastern. Is not he also just as clearly creating
value-making capital? The giant steamship, as the pick, is an article of
wealth, an instrument of production, and though the one may not be
completed for years, while the other is completed in a few minutes, each
day's work, in the one case as in the other, is as clearly a production
of wealth---an addition to capital. In the case of the steamship, as in
the case of the pick, it is not the last blow, any more than the first
blow, that creates the value of the finished product -the creation of
value is continuous, it immediately results from the exertion of labor.

We see this very clearly wherever the division of labor has made it
customary for different parts of the full process of production to be
carried on by different sets of producers---that is to say, wherever we
are in the habit of estimating the amount of value which the labor
expended in any preparatory stage of production has created. And a
moment's reflection will show that this is the case as to the vast
majority of products. Take a ship, a building, a jack-knife, a book, a
lady's thimble or a loaf of bread. They are finished products. But they
were not produced at one operation or by one set of producers. And this
being the case, we readily distinguish different points or stages in the
creation of the value which as completed articles they represent. When
we do not distinguish different parts in the final process of production
we do distinguish the value of the materials. The value of these
materials may often be again decomposed many times, exhibiting as many
clearly defined steps in the creation of the final value. At each of
these steps we habitually estimate a creation of value, an addition to
capital. The batch of bread which the baker is taking from the oven has
a certain value. But this is composed in part of the value of the flour
from which the dough was made. And this again is composed of the value
of the wheat, the value given by milling, etc. Iron in the form of pigs
is very far from being a completed product. It must yet pass through
several, or, perhaps, through many, stages of production before it
results in the finished articles that were the ultimate objects for
which the iron ore was extracted from the mine. Yet, is not pig iron
capital? And so the process of production is not really completed when
a. crop of cotton is gathered, nor yet when it is ginned and pressed;
nor yet when it arrives at Lowell or Manchester; nor yet when it is
converted into yarn; nor yet when it becomes cloth; but only when it is
finally placed in the hands of the consumer. Yet at each step in this
progress there is clearly enough a creation of value-an addition to
capital. Why, therefore, although we do not so habitually distinguish
and estimate it, is there not a creation of value-an addition to
capital-when the ground is plowed for the crop? Is it because it may
possibly be a bad season and the crop may fail? Evidently not; for a
like possibility of misadventure attends everyone of the many steps in
the production of the finished article. On the average a crop is sure to
come up, and so much plowing and sowing will on the average result in so
much cotton in the boll, as surely as so much spinning of cotton yarn
will result in so much cloth.

In short, as the payment of wages is always conditioned upon the
rendering of labor, the payment of wages in production, no matter how
long the process, never involves any advance of capital, or even
temporarily lessens capital. It may take a year, or even years, to build
a ship, but the creation of value of which the finished ship will be the
sum goes on day by day, and hour by hour, from the time the keel is laid
or even the ground is cleared. Nor by the payment of wages before the
ship is completed, does the master builder lessen either his capital or
the capital of the community, for the value of the partially completed
ship stands in place of the value paid out in wages. There is no advance
of capital in this payment of wages, for the labor of the workmen during
the week or month creates and renders to the builder more capital than
is paid back to them at the end of the week or month, as is shown by the
fact that if the builder were at any stage of the construction asked to
sell a partially completed ship he would expect a profit.

And so, when a Sutro or St.~Gothard tunnel or a Suez canal is cut, there
is no advance of capital. The tunnel or canal, as it is cut, becomes
capital as much as the money spent in cutting it-or, if you please, the
powder, drills, etc., used in the work, and the food, clothes, etc.,
used by the workmen-as is shown by the fact that the value of the
capital stock of the company is not lessened as capital in these forms
is gradually changed into capital in the form of tunnel or canal. On the
contrary, it probably, and on the average, increases as the work
progresses, just as the capital invested in a speedier mode of
production would on the average increase.

And this is obvious in agriculture also. That the creation of value does
not take place all at once when the crop is gathered, but step by step
during the whole process which the gathering of the crop concludes, and
that no payment of wages in the interim lessens the farmer's capital, is
tangible enough when land is sold or rented during the process of
production, as a plowed field will bring more than an unplowed field, or
a field that has been sown more than one merely plowed. It is tangible
enough when growing crops are sold, as is sometimes done, or where the
farmer does not harvest himself, but lets a contract to the owner of
harvesting machinery. It is tangible in the case of orchards and
vineyards which, though not yet in bearing, bring prices proportionate
to their age. It is tangible in the case of horses, cattle and sheep,
which increase in value as they grow toward maturity. And if not always
tangible between what may be called the usual exchange points in
production, this increase of value as surely takes place with every
exertion of labor. Hence, where labor is rendered before wages are paid,
the advance of capital is really made by labor, and is from the employed
to the employer, not from the employer to the employed.

``Yet,'' it may be said, ``in such cases as we have been considering
capital is required!'' Certainly; I do not dispute that. But it is not
required in order to make advances to labor. It is required for quite
another purpose. What that purpose is we may readily see.

When wages are paid in kind-that is to say, in wealth of the same
species as the labor produces; as, for instance, if I hire men to cut
wood, agreeing to give them as wages a portion of the wood they cut, a
method sometimes adopted by the owners or lessees of woodland, it is
evident that no capital is required for the payment of wages. Nor yet
when, for the sake of mutual convenience, arising from the fact that a
large quantity of wood can be more readily and more advantageously
exchanged than a number of small quantities, I agree to pay wages in
money, instead of wood, shall I need any capital, provided I can make
the exchange of the wood for money before the wages are due. It is only
when I cannot make such an exchange, or such an advantageous exchange as
I desire, until I accumulate a large quantity of wood that I shall need
capital. Nor even then shall I need capital if I can make a partial or
tentative exchange by borrowing on my wood. If I cannot, or do not
choose, either to sell the wood or to borrow upon it, and yet wish to go
ahead accumulating a large stock of wood, I shall need capital. But
manifestly, I need this capital, not for the payment of wages, but for
the accumulation of ·a stock of wood. Likewise in cutting a tunnel. If
the workmen were paid in tunnel (which, if convenient, might easily be
done by paying them in stock of the company), no capital for the payment
of wages would be required. It is only when the undertakers wish to
accumulate capital in the shape of a tunnel that they will need capital.
To recur to our first illustration: The broker to whom I sell my silver
cannot carry on his business without capital. But he does not need this
capital because he makes any advance of capital to me when he receives
my silver and hands me gold. He needs it because the nature of the
business requires the keeping of a certain amount of capital on hand, in
order that when a customer comes he may be prepared to make the exchange
the customer desires.

And so we shall find it in every branch of production. Capital has never
to be set aside for the payment of wages when the produce of the labor
for which the wages are paid is exchanged as soon as produced; it is
only required when this produce is stored up, or what is to the
individual the same thing, placed in the general current of exchanges
without being at once drawn against---that is, sold on credit. But the
capital thus required is not required for the payment of wages, nor for
advances to labor, as it is always represented in the produce of the
labor. It is never as an employer of labor that any producer needs
capital; when he does need capital, it is because he is not only an
employer of labor, but a merchant or speculator in, or an accumulator
of, the products of labor. This is generally the case with employers.

To recapitulate: The man who works for himself gets his wages in the
things he produces, as he produces them, and exchanges this value into
another form whenever he sells the produce. The man who works for
another for stipulated wages in money works under a contract of
exchange. He also creates his wages as he renders his labor, but he does
not get them except at stated times, in stated amounts, and in a
different form. In performing the labor he is advancing in exchange;
when he gets his wages the exchange is completed. During the time he is
earning the wages he is advancing capital to his employer, but at no
time, unless wages are paid before work is done, is the employer
advancing capital to him. Whether the employer who receives this produce
in exchange for the wages immediately re-exchanges it, or keeps it for
awhile, no more alters the character of the transaction than does the
final disposition of the product made by the ultimate receiver, who may,
perhaps, be in another quarter of the globe and at the end of a series
of exchanges numbering hundreds.

\newpage
